\scanpage{100}
Cette dernière condition est satisfaite pour le topos \inserted{$(S_s=)$} $\widehat\Delta$, car pour celui-ci
\[
\underbrace{\text{flasque}+\text{plat}\inserted{\uline{\unsure{non, projectif|lire projectif}}}}_{\Downarrow}\Longrightarrow\text{homotope à }0,
\]
donc \unsure{univ. flasque} et les $\Gamma^i(\Phi)$ univ. flasques $(i>0)$.
\margin{faux tel quel, à arranger}

Soit $P$ le topos ponctuel, et
\[
i_n:P\longrightarrow\widehat\Delta\ (S_s)
\]
le morphisme de topos associé au foncteur fibre
\[
i_n^*(K)=K_n=K(\Delta_n).
\]
Soit $k$ un anneau (constant)

\uline{Lemme} $\forall$ $k$-module $M$, on a un isom. can.
\[
i_{0*}(M)\simeq\Phi_*\otimes_{\mathbb Z}M\simeq\Phi_{*k}\otimes_k M
\]
où $\Phi_*\in\operatorname{Ob}(\widehat\Delta_{\mathrm{ab}}\simeq\operatorname{Hom}(\Delta^\circ,\mathrm{Ab.}))$ est défini par $\Phi_*(I)=\mathbb Z^I$, \unsure{donc|d'où} $\Phi_{*k}(I)\simeq k^I$.

En effet, on a par construction de l'image directe
\[
\begin{array}{ccc}
i_{0*}(M)(I)&\simeq M(i_0^*(I))&\simeq M^I\\
\Vert&\Vert&\\
\operatorname{Hom}(I,i_{0*}(M))&\operatorname{Hom}(\underbrace{i_0^*(I)}_{\operatorname{Hom}(\Delta_0,I)\simeq I},M)&
\end{array}
\]
\margin{OK.}
\deleted{\unsure{Soit le}}
\uline{Corollaire} Pour tout $k_*$-Module $A_*$, on a
\[
\operatorname{Hom}_{k_*}(A_*,\Phi_{*k}\otimes_k M)\simeq\operatorname{Hom}_k(A_0,M)
\]
en particulier on a un hom can. (\unsure{pour} $M=\unsure{C_0|A_0}$)
\[
\deleted{\unsure{Hom}}\quad A_*\longrightarrow\Phi_*\otimes_{\mathbb Z}A_0\simeq\Phi_{*k}\otimes_k A_0\ (=i_{0*}i_0^*A_*)
\]
Appliquons ceci au cas d'un monomorphisme
\[
k_*\xrightarrow{u} A_*=A_*^{1,0}
\]
provenant d'une $A_*^{**}\in\operatorname{Ob}\mathrm{DR}_{\unsure{\mathrm{sp}}}(S_s,k_*)$

\margin{99}
